% ----------------------------------------------------------
\chapter{Cronograma}
% ----------------------------------------------------------

Na tabela \ref{tab:cronograma} apresentada abaixo, as tarefas estão distribuídas ao longo de dois períodos distintos. As atividades enumeradas de 1 a 5 estão programadas para o período compreendido entre agosto e dezembro de 2023. Por outro lado, as tarefas de 6 a 11 estão agendadas para serem executadas no semestre que se estende de janeiro a junho de 2024. As tarefas são divididas do seguinte modo:

\begin{enumerate}
  \item \textbf{Definição do escopo e fluxos}: Esta fase envolve a definição clara do escopo do projeto e a identificação dos fluxos de trabalho necessários.
  \item \textbf{Implementação}: Nesta fase, a implementação real do projeto ocorre, com base no escopo e nos fluxos definidos anteriormente.
  \item \textbf{Documentação do código}: Aqui, a documentação detalhada do código desenvolvido é criada, facilitando a manutenção e a compreensão futura.
  \item \textbf{Deploy e testes}: O código é implantado em um ambiente e serão realizados testes extensivos para garantir que tudo está funcionando como esperado.
\end{enumerate}

\begin{table}[h]
\centering
\begin{tabular}{r*{11}{|c}|}
\cline{2-12} 
 & \multicolumn{11}{c|}{2023/2024} \\
\cline{2-12}
 & 1 & 2 & 3 & 4 & 5 & 6 & 7 & 8 & 9 & 10 & 11 \\
\cline{2-12}
\textbf{Definição} & X & X & X &   &   &   &   &   &   &   &   \\
\cline{2-12}
Atividade 1 & X & X & X &   &   &   &   &   &   &   &   \\
\cline{2-12}
\textbf{Implementação} &   &   & X & X & X & X & X & X & X &   &   \\
\cline{2-12}
Atividade 2 &   &   & X & X & X & X & X & X &   &   &   \\
\cline{2-12}
Atividade 3 &   &   &   &   &   &   & X & X & X &   &   \\
\cline{2-12}
Atividade 4 &   &   &   &   &   &   &   & X & X &   &   \\
\cline{2-12}
\textbf{Monografia} &   &   &   &   &   &   & X & X & X & X & X \\
\cline{2-12}
Escrita &   &   &   &   &   &   & X & X & X & X &   \\
\cline{2-12}
Defesa &   &   &   &   &   &   &   &   &   & X & X \\
\cline{2-12}
\end{tabular}
\caption{Cronograma das atividades do projeto}
\label{tab:cronograma}
\end{table}

% ----------------------------------------------------------
\chapter{Resultados Esperados}
% ----------------------------------------------------------

O objetivo do projeto Protótipo de um modelador de diagramas ER colaborativo é fornecer uma ferramenta capaz e acessível que atenda às demandas atuais de projetos de TI. Ao fornecer uma solução que simplifica de forma eficaz e colaborativa a modelagem de diagramas ER, esta iniciativa deve preencher uma lacuna significativa no mercado. A ferramenta em desenvolvimento promove a troca de ideias mais eficaz e a resolução de problemas em grupo.

Ademais, o projeto aspira fomentar um ambiente laboral mais produtivo, salientando a eficácia e reduzindo os custos vinculados ao desenvolvimento de projetos de banco de dados. Por meio da promoção da colaboração em tempo real, espera-se que a ferramenta atenue atrasos e retrabalhos, contribuindo para um ciclo de desenvolvimento mais ágil e dinâmico. Este empreendimento não apenas atende às demandas de um mundo digital e interconectado, mas também se posiciona como um catalisador para futuras inovações tecnológicas no campo da TI.

Ao final deste projeto, espera-se que o protótipo criado se torne uma ferramenta útil, ajudando a modelar bancos de dados e promovendo a cooperação e a eficiência no desenvolvimento de projetos de TI.